\documentclass[a4paper,11pt]{ltjsarticle}
\usepackage[margin=25mm]{geometry}
\usepackage{enumitem,listings}
\begin{document}
\begin{titlepage}
\noindent 上端の試料札。
\vfill
\noindent\begin{minipage}{.6\linewidth}
\small 折り板の保管手順。\\
青い面を上へ向ける。\\
窓の端に揃えて置く。
\end{minipage}
\vfill
\noindent 下端の試料札。
\end{titlepage}
\section{測定の準備}
\subsection{紙を折る}
\subsubsection*{窓の端}
最初に基準線を写す。
\begin{center}紙帯の中央\end{center}
\begin{quote}右端まで伸びる短い印を描く。\end{quote}
\begin{verse}
白い窓の影\\
青い窓の影

曇り窓の影
\end{verse}
\begin{itemize}[itemsep=3pt]
\item 二枚の紙
\begin{itemize}
\item 薄い紙
\item 厚い紙
\end{itemize}
\item 三枚の紙
\end{itemize}
\begin{lstlisting}[basicstyle=\small\ttfamily,numbers=left,caption={紙の枚数}]
white = 2
blue = 3
\end{lstlisting}
枚数を表の裏に記す。
\end{document}
